\documentclass[12pt]{article}
\usepackage[utf8]{inputenc}
\usepackage[T1]{fontenc}
\usepackage[margin=1in]{geometry}
\usepackage{setspace}
\usepackage{xcolor}
\usepackage[colorlinks=true, urlcolor=blue!60!black, linkcolor=blue!60!black]{hyperref}
\setlength{\parindent}{2em}
\setlength{\parskip}{1.25em}
\renewcommand{\baselinestretch}{1.0}
\title{\textit{Snyder v. Phelps}, 562 U.S. 443 (2011)}
\author{Case Brief}
\date{}
\begin{document}
\maketitle

% Case Brief: Snyder v. Phelps, 562 U.S. 443 (2011)
% Constitutional Law Briefing Book --- Chapter: Free Speech & Press

\section{Case Metadata}
\begin{itemize}
  \item \textbf{Case Name and Citation:} \textit{Snyder v.\ Phelps}, 562 U.S.\ 443 (2011) (No.\ 09-751)
  \item \textbf{Argued / Decided:} Argued October 6, 2010; decided March 2, 2011
  \item \textbf{Vote:} 8--1
  \item \textbf{Author:} Chief Justice John G.\ Roberts Jr.\ (joined by Justices Scalia, Kennedy, Thomas, Ginsburg, Breyer, Sotomayor, and Kagan)
  \item \textbf{Concurrences / Dissents:} Concurrence by Justice Stephen G.\ Breyer; dissent by Justice Samuel A.\ Alito Jr.
  \item \textbf{Chapter / Topic:} Free Speech \& Press --- Tort Liability, Matters of Public Concern, and Public Forum Picketing
  \item \textbf{Primary Source URLs:} \href{https://supreme.justia.com/cases/federal/us/562/443/}{Justia U.S.\ Supreme Court Center} $|$ \href{https://www.law.cornell.edu/supct/html/09-751.ZO.html}{Cornell LII}
\end{itemize}

\section{Facts and Procedural History}
Fred Phelps founded the Westboro Baptist Church in Topeka, Kansas, in 1955. Over more than two decades, Westboro parishioners picketed nearly 600 funerals to publicize their religious belief that God hates and punishes the United States for tolerating homosexuality, particularly in the military, as well as to condemn the Catholic Church over clergy abuse scandals.

Marine Lance Corporal Matthew Snyder was killed in Iraq in the line of duty, and his father, Albert Snyder, arranged a Catholic funeral in Westminster, Maryland, with time and location announced in local newspapers. On the day of the service, Phelps and six family-member parishioners picketed on public land near the Maryland State House, the U.S.\ Naval Academy, and Matthew Snyder's funeral. After notifying local authorities in advance and following police directions, the picketers stood behind a temporary fence on a 10-by-25-foot plot of public land adjacent to a public street, roughly 1,000 feet from the church and separated from it by several buildings. For about 30 minutes before the service, they sang hymns, recited Bible verses, and peacefully displayed signs reading, among other messages, ``God Hates the USA/Thank God for 9/11,'' ``America is Doomed,'' ``Thank God for IEDs,'' ``Thank God for Dead Soldiers,'' ``Pope in Hell,'' ``Priests Rape Boys,'' ``God Hates Fags,'' ``You're Going to Hell,'' and ``God Hates You.'' No picketer entered church property, yelled, used profanity, or committed violence. The funeral procession passed within 200 to 300 feet of the site; Albert Snyder saw only the tops of the signs while driving to the service and did not read their text until watching a television news broadcast that night. Weeks later, Westboro posted an online ``epic'' denouncing the Snyders directly, which the Supreme Court declined to consider because Snyder failed to raise it in his certiorari petition.

Snyder sued Phelps, two of his daughters, and the church in the U.S.\ District Court for the District of Maryland under diversity jurisdiction, asserting five Maryland state tort claims: defamation, publicity given to private life, intentional infliction of emotional distress (IIED), intrusion upon seclusion, and civil conspiracy. The District Court granted Westboro summary judgment on defamation and publicity given to private life, but a jury found for Snyder on IIED, intrusion upon seclusion, and civil conspiracy, awarding \$2.9 million in compensatory damages and \$8 million in punitive damages (remitted by the trial judge to \$2.1 million, for a \$5 million total). The Fourth Circuit reversed (580 F.3d 206 (4th Cir.\ 2009)), holding that Westboro's picketing was protected by the First Amendment as hyperbolic rhetoric on matters of public concern. The Supreme Court granted certiorari.

\section{Issue(s) Presented}
\begin{enumerate}
  \item Does the Free Speech Clause of the First Amendment shield peaceful picketers on public land near a private citizen's military funeral from state tort liability for intentional infliction of emotional distress when their signs address matters of public concern?
  \item Does the ``captive audience'' doctrine permit recovery for intrusion upon seclusion (and derivative civil conspiracy) where mourners at a funeral service could see only the tops of distant picket signs and the service itself was not disrupted?
\end{enumerate}

\section{Holding}
Yes to the first issue; no to the second. The First Amendment shields Westboro from tort liability for IIED, intrusion upon seclusion, and civil conspiracy arising from its peaceful picketing on matters of public concern in a traditional public forum. Judgment of the Fourth Circuit affirmed.

\section{Reasoning}
Writing for an eight-Justice majority, Chief Justice Roberts anchored the First Amendment defense in the distinction between speech on matters of public concern and speech of purely private significance:

\begin{itemize}
  \item \textbf{Public Concern Hierarchy:} Drawing on \href{https://supreme.justia.com/cases/federal/us/485/46/}{\textit{Hustler Magazine, Inc.\ v.\ Falwell}}, 485 U.S.\ 46 (1988), \textit{Connick v.\ Myers}, 461 U.S.\ 138 (1983), and \textit{Dun \& Bradstreet, Inc.\ v.\ Greenmoss Builders, Inc.}, 472 U.S.\ 749 (1985), Roberts reaffirmed that the Free Speech Clause operates as a defense to state tort suits and that speech on public issues occupies the ``highest rung of the hierarchy of First Amendment values.'' To determine whether speech addresses a matter of public concern, courts independently examine the ``content, form, and context'' of the speech across the whole record.
  \item \textbf{Content, Form, and Context of Westboro's Picketing:} Although Westboro's placards fell short of refined commentary, their \textit{content}---the moral and political conduct of the United States, homosexuality in the military, and Catholic clergy scandals---plainly addressed matters of broad public import, and their \textit{form} targeted a wide public audience. Even if a few signs (``You're Going to Hell,'' ``God Hates You'') could be read as directed at the Snyders, the ``overall thrust and dominant theme'' spoke to public issues. Nor did the funeral \textit{context} transform the speech into a private attack: Westboro had no pre-existing grudge or relationship with the Snyders and had protested on identical themes for 20 years before learning of Matthew Snyder.
  \item \textbf{Traditional Public Forum and ``Outrageousness'':} Westboro picketed peacefully on public land adjacent to a public street---``the archetype of a traditional public forum''---in full compliance with police instructions and 1,000 feet from the church. While content-neutral time, place, or manner rules (such as funeral buffer statutes enacted later by Maryland and 43 other states) raise separate questions not before the Court, tort liability here turned squarely on the content and viewpoint of Westboro's message. Echoing \textit{Hustler}, Roberts held that Maryland's ``outrageousness'' standard for IIED is too malleable and subjective, inviting juries to punish unpopular viewpoints and chilling public debate.
  \item \textbf{Rejection of Captive Audience, Intrusion, and Conspiracy Claims:} Emphasizing that the captive audience doctrine applies sparingly (primarily to protect the privacy of the home, as in \textit{Rowan v.\ Post Office Dept.} and \textit{Frisby v.\ Schultz}), Roberts declined to extend it where protesters remained out of sight, Snyder saw only sign tops, and the funeral was undisrupted. Because the First Amendment barred both IIED and intrusion upon seclusion, the derivative civil conspiracy claim necessarily failed.
\end{itemize}

\section{Concurring and Dissenting Opinions}
\begin{itemize}
  \item \textbf{Justice Breyer (concurring):} Joining the Court's opinion, Breyer stressed its narrowness, noting it did not address television broadcasts or Westboro's online ``epic.'' Rejecting an absolutist rule that public-concern speech always trumps state tort law, Breyer favored case-specific proportionality review: here, because Westboro complied with police directions and could not be seen or heard from the service, imposing IIED liability would punish public-concern speech without proportionately advancing Maryland's interest in preventing severe emotional harm.
  \item \textbf{Justice Alito (dissenting):} Writing alone, Alito argued that the First Amendment is not a license for a ``vicious verbal assault'' on private figures at a moment of acute grief. Emphasizing Westboro's pre-funeral press release, specific placards, and post-funeral online ``epic'' (which he faulted the majority for ignoring), Alito contended that Westboro deliberately attacked Matthew Snyder and his family to manufacture media publicity. Interspersing protected commentary on public issues should no more immunize a targeted personal attack than it immunizes defamation or fighting words, and a public street should not become a ``free-fire zone'' where private mourners lose the narrow, demanding protection of the IIED tort.
\end{itemize}

\section{Key Quotations}
\begin{enumerate}
  \item ``What Westboro said, in the whole context of how and where it chose to say it, is entitled to `special protection' under the First Amendment, and that protection cannot be overcome by a jury finding that the picketing was outrageous.'' --- Chief Justice Roberts, 562 U.S.\ at 458 (\href{https://supreme.justia.com/cases/federal/us/562/443/}{Justia}).
  \item ``Speech is powerful. It can stir people to action, move them to tears of both joy and sorrow, and---as it did here---inflict great pain. On the facts before us, we cannot react to that pain by punishing the speaker. As a Nation we have chosen a different course---to protect even hurtful speech on public issues to ensure that we do not stifle public debate.'' --- Chief Justice Roberts, 562 U.S.\ at 460--461 (\href{https://supreme.justia.com/cases/federal/us/562/443/}{Justia}).
  \item ``To uphold the application of state law in these circumstances would punish Westboro for seeking to communicate its views on matters of public concern without proportionately advancing the State's interest in protecting its citizens against severe emotional harm.'' --- Justice Breyer (concurring), 562 U.S.\ at 462--463 (\href{https://supreme.justia.com/cases/federal/us/562/443/}{Justia}).
  \item ``Our profound national commitment to free and open debate is not a license for the vicious verbal assault that occurred in this case. \dots\ In order to have a society in which public issues can be openly and vigorously debated, it is not necessary to allow the brutalization of innocent victims like petitioner.'' --- Justice Alito (dissenting), 562 U.S.\ at 463, 475 (\href{https://supreme.justia.com/cases/federal/us/562/443/}{Justia}).
\end{enumerate}

\section{Doctrinal Significance and Modern Legacy (Through 2026)}
\textit{Snyder v.\ Phelps} extended \href{https://supreme.justia.com/cases/federal/us/485/46/}{\textit{Hustler Magazine, Inc.\ v.\ Falwell}}'s First Amendment shield against IIED liability beyond public-figure targets to \textit{private-figure} plaintiffs whenever the defendant's speech, viewed in full context, addresses a matter of public concern in a public forum. It also cemented the Roberts Court's strong aversion to creating new categorical speech exceptions or allowing vague tort standards like ``outrageousness'' to function as backdoor viewpoint censorship.

Through 2026, \textit{Snyder}'s core anti-censorship ethos has proven remarkably durable across the Court's First Amendment jurisprudence:
\begin{itemize}
  \item \textbf{Buffer Zones and Time, Place, or Manner Limits:} While \textit{Snyder} pointed to content-neutral funeral-picketing statutes as a constitutionally distinct alternative to IIED verdicts, the Court later applied rigorous narrow-tailoring review to public-sidewalk buffer zones in \href{https://supreme.justia.com/cases/federal/us/573/464/}{\textit{McCullen v.\ Coakley}}, 573 U.S.\ 464 (2014) (invalidating a 35-foot abortion-clinic buffer zone), confirming that even content-neutral spatial restrictions in traditional public forums face demanding scrutiny.
  \item \textbf{``Breathing Space'' Against Tort and Criminal Liability:} The Court has repeatedly reinforced \textit{Snyder}'s protection of offensive, caustic, or emotionally distressing expression outside narrow historic exceptions---striking down the Stolen Valor Act in \href{https://supreme.justia.com/cases/federal/us/567/709/}{\textit{United States v.\ Alvarez}}, 567 U.S.\ 709 (2012), protecting vulgar off-campus student criticism in \href{https://www.supremecourt.gov/opinions/20pdf/20-255_g3bi.pdf}{\textit{Mahanoy Area School District v.\ B.L.}}, 594 U.S.\ 180 (2021), requiring a subjective recklessness \textit{mens rea} for ``true threats'' prosecutions (including stalking-by-speech causing severe emotional distress) in \href{https://supreme.justia.com/cases/federal/us/600/22-138/}{\textit{Counterman v.\ Colorado}}, 600 U.S.\ 66 (2023), barring government coercion to suppress disfavored advocacy in \href{https://supreme.justia.com/cases/federal/us/602/22-842/}{\textit{NRA v.\ Vullo}}, 602 U.S.\ 175 (2024), and reaffirming in \href{https://supreme.justia.com/cases/federal/us/607/24-539/}{\textit{Chiles v.\ Salazar}}, 607 U.S.\ \_\_\_ (2026), that states cannot evade strict scrutiny for content-based speech restrictions by invoking broad tort or regulatory analogies outside narrow historical traditions.
  \item \textbf{Unresolved Digital Frontier:} Because \textit{Snyder} expressly bypassed Westboro's online ``epic,'' lower courts through 2026 continue to wrestle with applying \textit{Snyder}'s ``public concern'' and ``context'' framework to targeted online harassment, doxxing, and social-media campaigns against private individuals.
\end{itemize}

\section{Pedagogical Critiques and Exam Hypotheticals}
\textbf{Pedagogical Critiques:}
\begin{itemize}
  \item \textbf{The ``Pre-Existing Grudge'' Paradox:} As Alito sharply observed, Roberts's distinction between a contrived personal vendetta and Westboro's impersonal publicity strategy yields a perverse moral inversion: deliberately inflicting severe emotional trauma on a grieving stranger to attract television cameras receives \textit{greater} constitutional immunity than lashing out in a private feud.
  \item \textbf{Indeterminacy of ``Public Concern'':} Courts have struggled to police the boundary between public and private concern without making ad hoc content judgments. If adding a few generic political slogans to targeted personal abuse shifts the ``dominant theme'' to public concern, \textit{Snyder} risks creating an easy blueprint for immunizing private harassment.
\end{itemize}

\textbf{Exam Hypothetical:} After a high-profile fatal car crash involving a drunk driver, an anti-drunk-driving activist group pickets outside the home of the deceased driver's grieving parents on a public sidewalk for three consecutive nights, carrying signs reading ``Drunk Drivers Are Murderers,'' ``Your Son Deserved to Die,'' and ``Weak DUI Laws Kill.'' The group also posts the parents' private home address and phone number on social media urging followers to ``make them feel the pain.'' The parents sue for IIED and intrusion upon seclusion.
\begin{itemize}
  \item \textbf{Activists' Best Argument (Applying \textit{Snyder}):} Under \textit{Snyder}, DUI enforcement and highway safety are classic matters of public concern, the picketing occurred on a traditional public sidewalk, the group had no prior private dispute with the parents, and a jury cannot impose IIED liability merely because the signs are cruel or ``outrageous.''
  \item \textbf{Parents' Best Argument (Opposite-Side Contrast):} Distinguish \textit{Snyder} on three decisive grounds. First, unlike a distant 1,000-foot church-adjacent plot, targeted picketing directly outside a private residence implicates the home's heightened privacy and the captive audience doctrine (\textit{Frisby v.\ Schultz}), supporting both intrusion upon seclusion and content-neutral residential picketing limits. Second, the online doxxing campaign---posting private contact details and urging followers to ``make them feel the pain''---resembles Westboro's unadjudicated online ``epic'' and goes beyond public debate into targeted private harassment (and potentially reckless ``true threats'' under \textit{Counterman v.\ Colorado}). Third, under Breyer's proportionality concurrence and Alito's dissent, the combination of residential targeting and doxxing directly invades private seclusion without advancing public debate.
\end{itemize}

\end{document}
